\documentclass[11pt]{article}
\usepackage[margin=0.85in]{geometry}
\usepackage{booktabs}
\usepackage{array}
\usepackage{longtable}
\usepackage{float}
\usepackage{hyperref}

\title{Project 1A LSQ Fund Update}
\author{Prepared for Professor Phelim Boyle}
\date{Updated through Jun 2026}

\begin{document}
\maketitle

\section*{Data \& Method}
The update uses 114 monthly LSQ observations from January 2017 through Jun 2026.
The original paper used 109 observations ending in January 2026; the added LSQ returns are Feb 2026 = 1.88\%, Mar 2026 = 3.68\%, Apr 2026 = 2.95\%, May 2026 = 7.64\%, Jun 2026 = 6.39\%.
The LSQ analysis uses the baseline net-of-fees return series from the Spartan monthly performance PDF.

For the benchmark, I used the Yahoo Finance \texttt{\string^SP500TR} S\&P 500 Total Return Index.
This is the appropriate S\&P 500 series because the paper states that the benchmark returns are calculated with dividends reinvested; a price-only S\&P 500 index would omit dividends and understate the benchmark return.

For 2026-06, the market comparators are calculated from Yahoo Finance daily adjusted closes at each month-end. The S\&P 500 TR return is -0.952323\%, and the NVIDIA return is -5.123049\%.

All return inputs are monthly.
Sample standard deviation divides by $N-1$.
Skewness and excess kurtosis use the adjusted sample formulas.
Sharpe ratios use a 0\% risk-free rate.
Sortino ratios use the paper's downside-deviation convention.
Omega is computed as $\sum \max(r_t-L,0) / \sum \max(L-r_t,0)$.
Alpha and beta are estimated from the monthly regression of LSQ returns on S\&P 500 Total Return Index returns.

\clearpage
\section*{Updated Table 2}
\setcounter{table}{1}
\begin{table}[H]
\centering
\caption{Comparison of LSQ returns with S\&P 500 returns}
\label{tab:updated-table-2}
\begin{tabular}{lcc}
\toprule
\textbf{Metric} & \textbf{LSQ} & \textbf{S\&P 500 Index} \\
\midrule
\multicolumn{3}{l}{\emph{Statistical properties of monthly returns}} \\
\midrule
Average (arithmetic) return & 3.52 & 1.30 \\
Average (geometric) return & 3.45 & 1.19 \\
Percent Positive Months & 95.65 & 69.57 \\
Percent Negative Months & 4.35 & 30.43 \\
Average Positive Months & 3.72 & 3.62 \\
Average Negative Months & -0.87 & -4.01 \\
Best Month & 21.97 & 12.82 \\
Worst Month & -2.20 & -12.35 \\
\midrule
\multicolumn{3}{l}{\emph{First four moments of distribution of monthly returns}} \\
\midrule
Mean & 3.52 & 1.30 \\
Standard deviation & 3.98 & 4.51 \\
Skewness & 2.64 & -0.43 \\
Excess kurtosis & 8.54 & 0.54 \\
\midrule
\multicolumn{3}{l}{\emph{Risk adjusted performance metrics}} \\
\midrule
Sharpe Ratio & 0.88 & 0.29 \\
Sortino Ratio & 14.58 & 0.47 \\
Omega Ratio & 94.26 & 2.06 \\
Alpha & 3.55 & 0 \\
Beta & -0.03 & 1 \\
\bottomrule
\end{tabular}

\vspace{0.2cm}
\begin{minipage}{0.92\linewidth}
\footnotesize
Note: Returns are monthly percentages except ratios, skewness, excess kurtosis, and beta.
The LSQ column uses the baseline net-of-fees series.
The S\&P 500 benchmark is the total return index, consistent with the paper's dividend-reinvestment convention.
\end{minipage}
\end{table}


\clearpage
\section*{Updated Table 4}
\setcounter{table}{3}
\begin{table}[H]
\centering
\caption{Monthly Return Statistics: LSQ Fund and NVIDIA}
\label{tab:updated-table-4}
\small
\begin{tabular}{lcc}
\toprule
\textbf{Statistic} & \textbf{LSQ Fund} & \textbf{NVIDIA} \\
& \emph{Jan 2017 to Jul 2026} & \emph{Jan 2017 to Jul 2026} \\
& \emph{(N = 115)} & \emph{(N = 115)} \\
\midrule
\multicolumn{3}{l}{\emph{Panel A: Return Summary}} \\
\midrule
Arithmetic Mean Return & 3.52 & 4.72 \\
Geometric Mean Return & 3.45 & 3.84 \\
Percent Positive Months & 95.65 & 64.35 \\
Percent Negative Months & 4.35 & 35.65 \\
Average Positive Month & 3.72 & 12.21 \\
Average Negative Month & -0.87 & -8.79 \\
Best Month & 21.97 & 38.54 \\
Worst Month & -2.20 & -32.03 \\
\midrule
\multicolumn{3}{l}{\emph{Panel B: First Four Moments}} \\
\midrule
Mean ($\mu$) & 3.52 & 4.72 \\
Standard Deviation ($\sigma$) & 3.98 & 13.45 \\
Skewness & 2.64 & -0.02 \\
Excess Kurtosis & 8.54 & 0.15 \\
\midrule
\multicolumn{3}{l}{\emph{Panel C: Risk-Adjusted Ratios (Monthly, risk free rate = 0)}} \\
\midrule
Sharpe Ratio & 0.88 & 0.35 \\
Sortino Ratio & 14.58 & 0.67 \\
Omega Ratio & 94.26 & 2.51 \\
\bottomrule
\end{tabular}

\vspace{0.15cm}
\begin{minipage}{0.92\linewidth}
\footnotesize
Sortino uses Estrada (1999) downside deviation:
$\sigma_d=\sqrt{\frac{1}{N}\sum_{r_t<0}r_t^2}$.
Omega $= \sum \max(r_t,0) / \sum \max(-r_t,0)$.
\end{minipage}
\end{table}


\clearpage
\section*{Summary of Changes}
\subsection*{Table 2: Paper Values vs Updated Values}
\small
\begin{center}
\begin{tabular}{lrrr}
\toprule
Metric & Paper & Updated & Change \\
\midrule
LSQ arithmetic monthly return & 3.48\% & 3.52\% & +0.04 pp \\
LSQ geometric monthly return & 3.40\% & 3.45\% & +0.05 pp \\
LSQ percent positive months & 95.40\% & 95.61\% & +0.21 pp \\
LSQ percent negative months & 4.60\% & 4.39\% & -0.21 pp \\
LSQ sample standard deviation & 4.06\% & 4.00\% & -0.06 pp \\
LSQ Sharpe ratio & 0.86 & 0.88 & +0.02 \\
LSQ Sortino ratio & 13.96 & 14.53 & +0.57 \\
LSQ Omega ratio & 88.28 & 93.48 & +5.20 \\
Alpha, monthly & 3.52\% & 3.55\% & +0.03 pp \\
Beta & -0.04 & -0.03 & +0.01 \\
\bottomrule
\end{tabular}
\end{center}
\normalsize

\subsection*{Table 4: Paper LSQ Column vs Updated LSQ Column}
\small
\begin{center}
\begin{tabular}{lrrr}
\toprule
Statistic & Paper LSQ & Updated LSQ & Change \\
\midrule
Arithmetic Mean Return & 3.47 & 3.52 & +0.05 \\
Geometric Mean Return & 3.40 & 3.45 & +0.05 \\
Percent Positive Months & 95.41 & 95.61 & +0.20 \\
Percent Negative Months & 4.59 & 4.39 & -0.20 \\
Average Positive Month & 3.68 & 3.72 & +0.04 \\
Average Negative Month & -0.87 & -0.87 & +0.00 \\
Standard Deviation & 4.06 & 4.00 & -0.06 \\
Sharpe Ratio & 0.86 & 0.88 & +0.02 \\
Sortino Ratio & 13.96 & 14.53 & +0.57 \\
Omega Ratio & 88.28 & 93.48 & +5.20 \\
\bottomrule
\end{tabular}
\end{center}
\normalsize

\subsection*{Interpretation}
The updated LSQ statistics move only moderately because 5 months are being added to a 109-month sample.
The added sample contains 5 positive months and 0 negative months, leaving 109 positive months and 5 negative months in total.
The average monthly return rises relative to the paper after incorporating returns through Jun 2026.
The worst monthly return remains -2.20\% because none of the added months is lower.

The risk-adjusted metrics remain very strong.
The Sharpe ratio is 0.88, the Sortino ratio is 14.53, and the zero-threshold Omega ratio is 93.48.
The updated monthly alpha is 3.55\% and beta is -0.03, so the evidence still supports the paper's conclusion that LSQ's returns are not explained by broad equity-market exposure.

\clearpage
\section*{Additional Analysis: Serial Correlation}
Since the paper already identifies serial correlation as an important feature of the LSQ return series, I implemented a serial-correlation diagnostic table and a Newey-West adjusted inference table.
Newey-West inference adjusts the standard errors for autocorrelation and heteroskedasticity in the regression residuals.
This gives a more appropriate t-statistic when monthly returns are not independent over time.

\subsection*{Serial-Correlation Diagnostics}
\begin{table}[H]
\centering
\caption{Serial-Correlation Diagnostics for Monthly Returns}
\label{tab:serial-correlation}
\small
\begin{tabular}{lrrrr}
\toprule
Series & $\rho_1$ & LB p(1) & LB p(6) & LB p(12) \\
\midrule
LSQ Fund & 0.35 & 0.0001 & 0.0001 & 0.0001 \\
S\&P 500 TR & -0.14 & 0.1255 & 0.4456 & 0.3139 \\
LSQ minus S\&P 500 TR & 0.01 & 0.8867 & 0.9770 & 0.8757 \\
\bottomrule
\end{tabular}

\vspace{0.15cm}
\begin{minipage}{0.86\linewidth}
\footnotesize
$\rho_1$ is the first-order monthly autocorrelation.
LB p($k$) is the Ljung-Box p-value through lag $k$.
\end{minipage}
\end{table}

\subsection*{Newey-West Adjusted Inference}
\begin{table}[H]
\centering
\caption{Conventional vs Newey-West Inference}
\label{tab:newey-west}
\small
\begin{tabular}{lrrrrr}
\toprule
Statistic & Estimate & OLS SE & OLS t-stat & Newey-West SE & Newey-West t-stat \\
\midrule
LSQ monthly mean & 3.52\% & 0.37\% & 9.39 & 0.53\% & 6.66 \\
Monthly alpha & 3.55\% & 0.39\% & 9.07 & 0.58\% & 6.08 \\
Beta & -0.03 & 0.08 & -0.30 & 0.11 & -0.24 \\
\bottomrule
\end{tabular}

\vspace{0.15cm}
\begin{minipage}{0.92\linewidth}
\footnotesize
Newey-West standard errors use 4 monthly lags, based on the rule $\lfloor 4(T/100)^{2/9}\rfloor$ with $T=114$.
The alpha-beta regression is $r_{\mathrm{LSQ},t}=\alpha+\beta r_{\mathrm{S\&P500TR},t}+\varepsilon_t$.
\end{minipage}
\end{table}

The Newey-West adjustment increases the standard errors, which is expected when there is positive serial correlation.
Even after this adjustment, the LSQ monthly mean and alpha remain economically large and statistically strong.
The beta remains close to zero and statistically insignificant.

\subsection*{Additional Suggestions}
I think the most logical next extension would be to apply the same dependence-aware approach to the paper's nonlinear performance ratios.
We could produce confidence intervals for the Sharpe, Sortino, and Omega ratios without assuming independent monthly observations.
This should directly address the serial-correlation issue and preserve our focus on asymmetric performance.

\end{document}
